\begin{definition}[Stiefel 多様体]
$n,p$ を自然数で $p\le n$ とする．
実数体上の Stiefel 多様体 $\mathrm{St}(p,n)$ を
\[
\mathrm{St}(p,n)
=
\{\, W\in\mathbb{R}^{n\times p}\mid W^\top W = I_p \,\}
\]
によって定義する．
\end{definition}

\noindent
すなわち $\mathrm{St}(p,n)$ は，
$\mathbb{R}^n$ における $p$ 本の互いに直交し，
かつ単位長を持つ列ベクトルの組全体からなる集合である．
各行列 $W=[w_1,\dots,w_p]$ の列ベクトルは
\[
\langle w_i , w_j \rangle = \delta_{ij}
\]
を満たし，ここで $\delta_{ij}$ は Kronecker のデルタである．

\begin{remark}
幾何学的には，$\mathrm{St}(p,n)$ の各点 $W$ は
$\mathbb{R}^n$ 内の $p$ 次元部分空間の
\emph{正規直交基底}を表す．
ただし基底の取り方は一意ではなく，
任意の直交行列 $Q\in O(p)$ に対して
$W$ と $WQ$ は同じ部分空間を張る．
この同一視を行うと Grassmann 多様体
\[
\mathrm{Gr}(p,n)
\]
が得られる．
\end{remark}

\begin{remark}
$\mathrm{St}(p,n)$ は
$\mathbb{R}^{n\times p}$ の部分集合であり，
制約条件 $W^\top W=I_p$ によって定まる
滑らかな多様体となる．
その次元は
\[
\dim \mathrm{St}(p,n)=np-\frac{p(p+1)}{2}
\]
で与えられる．
\end{remark}

\begin{example}
$p=1$ のとき，
\[
\mathrm{St}(1,n)=\{w\in\mathbb{R}^n\mid \|w\|=1\}
\]
となり，これは $n$ 次元ユークリッド空間における単位球面 $S^{n-1}$ に一致する．
\end{example}
